\subsection{对正则序列取商}

命题 4.—— 设 A 为 Noether 环，J 为由 A-正则序列 x 生成的 A 的理想，而 Ω 为有限生成的 A-模。

\begin{enumerate}
\def\labelenumi{\alph{enumi})}
\item
  若 A-模 \(\Omega\) 是对偶化的，则序列 \(\mathbf{x}\) 是 \(\Omega\)-正则的，且 A/J-模 \(\Omega/\mathrm{J}\Omega\) 是对偶化的；
\item
  若 A/J-模 \(\Omega/\mathrm{J}\Omega\) 是对偶化的，J 含于 \(\Lambda\) 的根中，且序列 x 是 \(\Omega\)-正则的，则 A-模 \(\Omega\) 是对偶化的。
\end{enumerate}

对序列 x 的长度作归纳，可化归为它只含单个元素 x 的情形。设 A-模 \(\Omega\) 是对偶化的。对 A 的每个包含 x 的极大理想 m，有 \(\dim(\mathrm{A}_{\mathfrak{m}}/x\mathrm{A}_{\mathfrak{m}})=\dim(\mathrm{A}_{\mathfrak{m}})-1\)（VIII, § 3, 第 1 节，推论 2），从而 \(\mathrm{Hom}_{\mathrm{A}_{\mathfrak{m}}}\left(\mathrm{A}_{\mathfrak{m}} / x\mathrm{A}_{\mathfrak{m}}, \Omega_{\mathfrak{m}}\right) = 0\)（第 1 节，命题 3, a)）。这蕴含 \(\mathrm{Hom}_{\mathrm{A}}(\mathrm{A} / x\mathrm{A}, \Omega) = 0\)，于是乘法映射 \(x_{\Omega}\) 是单射的。因此，为证明命题，可设乘法映射 \(x_{\Omega}\) 是单射的。

令 \(\overline{\mathsf{A}}\) 为环 A/xA；设 \(\mathfrak{m}\) 为 A 的一个包含 \(x\) 的极大理想，并用 \(\overline{\mathfrak{m}}\) 表示它在 \(\overline{\mathsf{A}}\) 中的像。A-模 A/\(\mathfrak{m}\) 被 \(x\) 零化，且等同于 \(\overline{\mathsf{A}}/\overline{\mathfrak{m}}\)；于是对每个整数 \(i \geqslant 1\)，有一个同构 \(\operatorname{Ext}_{\mathsf{A}}^{i}(\mathsf{A}/\mathfrak{m}, \Omega) \longrightarrow \operatorname{Ext}_{\overline{\mathsf{A}}}^{i-1}(\overline{\mathsf{A}}/\overline{\mathfrak{m}}, \Omega/x\Omega)\)（\(\S\) 3, 第 4 节，命题 7）。我们有

\[
\mathrm{ht} (\overline {{\mathfrak {m}}}) = \dim (\overline {{\mathrm{A}}} _ {\overline {{\mathfrak {m}}}}) = \dim (\mathrm{A} _ {\mathfrak {m}} / x \mathrm{A} _ {\mathfrak {m}}) = \dim (\mathrm{A} _ {\mathfrak {m}}) - 1 = \mathrm{ht} (\mathfrak {m}) - 1
\]

（VIII, § 3, 第 1 节，推论 2, a)）。现在，\(\overline{\mathrm{A}}\) 的极大理想正是这样的理想 \(\overline{\mathfrak{m}}\)，其中 \(\mathfrak{m}\) 是 A 的包含 \(x\) 的极大理想；此外，若 \(x\) 属于 A 的根，则 A 的每个极大理想都包含 \(x\)。命题由此得证。

推论 1.—— 设 A 为 Noether 整环。每个对偶化 A-模都是无扭的且秩为 1。

设 \(\Omega\) 为对偶化 A-模；由命题 4，它是无扭的。设 K 为 A 的分式域；K-向量空间 \(K \otimes_{\Lambda} \Omega\) 是对偶化的（\(n^{\circ}\) 1, 命题 2），故维数为 1。

推论 2.—— 设 A 为局部 Macaulay 环，Ω 为有限生成的 A-模，而 x 为 m\(_{A}\) 中元素的一个极大割序列，它生成理想 J。下列条件等价：

\begin{enumerate}
\def\labelenumi{(\roman{enumi})}
\item
  A-模 Ω 是对偶化的；
\item
  A-模 \(\Omega\) 是 Macaulay 模，维数等于 \(\dim(A)\)，且 \(\Omega/J\Omega\) 是局部 Artin 环 \(A/J\) 上的不可分解内射模；
\item
  序列 x 是 Ω-正则的，且 Ω/JΩ 是局部 Artin 环 A/J 上的不可分解内射模；
\item
  序列 \(\mathbf{x}\) 是 \(\Omega\)-正则的，有 \(\mathrm{long}_{\mathrm{A}}(\Omega/\mathrm{J}\Omega) = \mathrm{long}_{\mathrm{A}}(\Lambda/\mathrm{J})\)，且 \(\kappa_{\mathrm{A}}\)-向量空间 \(\mathrm{Hom}_{\Lambda}(\kappa_{\Lambda}, \Omega/\mathrm{J}\Omega)\) 的维数为 1。
\end{enumerate}

(i)⇒(ii)：若 Ω 是对偶化的，则它是 Macaulay 模且维数为 dim(A)（第 1 节，命题 1）。由于 A 是 Macaulay 环，序列 x 是 A-正则的；由命题 4，A/J-模 Ω/JΩ 是对偶化的，从而是 Matlis A/J-模（第 1 节，例 1）。

(ii)⇒(iii)：在假设 (ii) 下，有 \(\dim(\Omega)=\dim(A)\) 与 \(\dim(\Omega/J\Omega)=\dim(A/J)=0\)，故序列 x 对 \(\Omega\) 是割的，从而是 \(\Omega\)-正则的（§ 2, 第 3 节，定理 1）。

(iii)⇒(i)：在 (iii) 的假设下，A/J-模 Ω/JΩ 是 Matlis A/J-模，从而是对偶化的（第 1 节，例 1）；由命题 4，A-模 Ω 是对偶化的。

\begin{enumerate}
\def\labelenumi{(\roman{enumi})}
\setcounter{enumi}{2}
\tightlist
\item
  \(\Leftrightarrow\)（iv）：这由 § 8, 第 3 节的注得出。
\end{enumerate}

